\documentclass{article}
\usepackage{amsmath,amssymb,amsthm}
\usepackage{fullpage}
\usepackage{hyperref}
\newtheorem{theorem}{Theorem}
\newtheorem{lemma}{Lemma}
\begin{document}

\section*{Primal–Dual Hybrid Gradient (PDHG) Algorithm}

\section*{Mathematical Formulation}
We consider a convex–concave saddle‑point problem
\[
\min_{x\in\mathbb R^{n}} \; \max_{y\in\mathbb R^{m}} \;
\mathcal L(x,y)\;:=\; f(x)+\langle Kx ,\,y\rangle - g^{*}(y),
\tag{1}
\]
where  

\begin{itemize}
\item $f:\mathbb R^{n}\rightarrow\mathbb R$ is convex and $L_{f}$‑smooth,
\item $g^{*}:\mathbb R^{m}\rightarrow\mathbb R$ is convex (the Fenchel conjugate of a convex $g$),
\item $K\in\mathbb R^{m\times n}$ is a linear operator with operator norm $\|K\|$.
\end{itemize}

Let $\tau>0$ (primal step) and $\sigma>0$ (dual step) satisfy the classic PDHG condition
\[
\tau\sigma\|K\|^{2}<1 .
\tag{2}
\]

\bigskip
\noindent\textbf{Iterates.} Starting from $(x_{0},y_{0})$, for $k=0,1,2,\dots$ compute
\[
\boxed{
\begin{aligned}
y_{k+1}&= y_{k}+ \sigma\bigl(Kx_{k}-\nabla g^{*}(y_{k})\bigr), \tag{3a}\\[4pt]
x_{k+1}&= \operatorname{prox}_{\tau f}\!\Bigl(x_{k}-\tau K^{\top}y_{k+1}\Bigr)
      \;=\; \arg\min_{x}\Bigl\{ f(x)+\frac{1}{2\tau}\|x- (x_{k}-\tau K^{\top}y_{k+1})\|^{2}\Bigr\}. \tag{3b}
\end{aligned}}
\]
If $f$ is $L_{f}$‑smooth, \eqref{3b} can be written explicitly as a gradient step
\[
x_{k+1}=x_{k}-\tau\bigl(\nabla f(x_{k})+K^{\top}y_{k+1}\bigr).
\tag{4}
\]

The algorithm stops after a predetermined number $T$ of iterations or when a
primal–dual gap tolerance is met.

\section*{Pseudocode}
\begin{verbatim}
Input:   step sizes τ>0, σ>0 with τσ‖K‖²<1
         initial points x0∈R^n, y0∈R^m
         maxIter = T
for k = 0,…,T-1 do
    # Dual gradient ascent
    y_{k+1} = y_k + σ ( K x_k - ∇g* (y_k) )
    # Primal proximal (or gradient) step
    x_{k+1} = prox_{τ f} ( x_k - τ Kᵀ y_{k+1} )
end for
Output: (x_T , y_T)   (or averaged iterates)
\end{verbatim}

\section*{Dry Run Example}
We illustrate the scalar case $n=m=1$, $K=1$, $g^{*}(y)=0$ (hence $\nabla g^{*}(y)=0$).  
Take
\[
f(x)=(x-3)^{2},\qquad \nabla f(x)=2(x-3).
\]
Choose step sizes $\tau=\sigma=0.1$ (they satisfy $\tau\sigma K^{2}=0.01<1$).  
Initialize $x_{0}=0$, $y_{0}=0$.

\bigskip
\noindent\textbf{Iteration 0}
\[
\begin{aligned}
y_{1}&=y_{0}+0.1\,(Kx_{0})=0+0.1\cdot0=0,\\
x_{1}&=x_{0}-0.1\bigl(\nabla f(x_{0})+K^{\top}y_{1}\bigr)
      =0-0.1\bigl(2(0-3)+0\bigr)=0-0.1(-6)=0.6 .
\end{aligned}
\]
Objective $f(x_{1})=(0.6-3)^{2}=5.76$.

\bigskip
\noindent\textbf{Iteration 1}
\[
\begin{aligned}
y_{2}&=y_{1}+0.1\,(Kx_{1})=0+0.1\cdot0.6=0.06,\\
x_{2}&=x_{1}-0.1\bigl(2(x_{1}-3)+y_{2}\bigr)
      =0.6-0.1\bigl(2(0.6-3)+0.06\bigr)\\
     &=0.6-0.1\bigl(-4.8+0.06\bigr)=0.6-0.1(-4.74)=1.074 .
\end{aligned}
\]
Objective $f(x_{2})=(1.074-3)^{2}=3.704.\,$

\bigskip
\noindent\textbf{Iteration 2}
\[
\begin{aligned}
y_{3}&=y_{2}+0.1\,x_{2}=0.06+0.1\cdot1.074=0.1674,\\
x_{3}&=x_{2}-0.1\bigl(2(x_{2}-3)+y_{3}\bigr)\\
     &=1.074-0.1\bigl(2(1.074-3)+0.1674\bigr)\\
     &=1.074-0.1\bigl(-3.852+0.1674\bigr)=1.074-0.1(-3.6846)=1.44246 .
\end{aligned}
\]
Objective $f(x_{3})=(1.44246-3)^{2}=2.424.$

The iterates move steadily toward the minimizer $x^{\star}=3$ while the dual variable $y$ grows, illustrating the primal–dual coupling.

\newpage
\section*{Assumptions}
\begin{enumerate}
\item \textbf{Convexity and smoothness of $f$.}  
   $f$ is convex and differentiable with $L_{f}$‑Lipschitz gradient:
   \[
   \|\nabla f(x)-\nabla f(x')\|\le L_{f}\|x-x'\|,\qquad\forall x,x'.
   \tag{A1}
   \]
\item \textbf{Convexity of $g^{*}$.}  
   $g^{*}$ is convex (possibly nonsmooth). Its subgradient $\partial g^{*}(y)$ exists for all $y$.
   \tag{A2}
\item \textbf{Linear operator $K$.}  
   $K$ is a known matrix with spectral norm $\|K\|$.
   \tag{A3}
\item \textbf{Step‑size condition.}  
   The primal and dual steps satisfy
   \[
   \tau\sigma\|K\|^{2}<1 .
   \tag{A4}
   \]
\item \textbf{Existence of a saddle point.}  
   There exists $(x^{\star},y^{\star})$ such that
   \[
   \mathcal L(x^{\star},y)\le\mathcal L(x^{\star},y^{\star})\le\mathcal L(x,y^{\star}),
   \qquad\forall x,y .
   \tag{A5}
   \]
\end{enumerate}

\section*{Main Convergence Theorem}
\begin{theorem}[Ergodic $O(1/T)$ primal–dual gap]\label{thm:rate}
Under assumptions (A1)–(A5) and the step‑size condition (A4), let
\[
\bar x_{T}:=\frac{1}{T}\sum_{k=0}^{T-1}x_{k},\qquad
\bar y_{T}:=\frac{1}{T}\sum_{k=0}^{T-1}y_{k+1}.
\]
Then the averaged iterates satisfy
\[
\boxed{
\mathcal G_{T}:=
\max_{y}\mathcal L(\bar x_{T},y)-\min_{x}\mathcal L(x,\bar y_{T})
\;\le\;
\frac{1}{2\tau T}\|x_{0}-x^{\star}\|^{2}
+\frac{1}{2\sigma T}\|y_{0}-y^{\star}\|^{2}.
}
\tag{5}
\]
Consequently $\mathcal G_{T}=O(1/T)$.
\end{theorem}

\section*{Theoretical Properties}
\begin{itemize}
\item \textbf{Stability.} Condition \eqref{2} guarantees that the Lyapunov
  quantity
  \[
  \Phi_{k}:=\frac{1}{2\tau}\|x_{k}-x^{\star}\|^{2}
           +\frac{1}{2\sigma}\|y_{k}-y^{\star}\|^{2}
  \]
  is non‑increasing up to a telescoping term (see Lemma~\ref{lem:descent}).
\item \textbf{Hyper‑parameter sensitivity.} The bound \eqref{5} is
  monotone decreasing in $\tau$ and $\sigma$; however they cannot be increased
  arbitrarily because of \eqref{2}. The optimal trade‑off is attained when
  $\tau\sigma\|K\|^{2}$ is close to $1$.
\item \textbf{Comparison to baselines.}
  \begin{itemize}
    \item For pure gradient descent on the primal only, the rate is $O(1/\sqrt{T})$
      for nonsmooth $g^{*}$, while PDHG attains $O(1/T)$ without requiring
      strong convexity.
    \item When $g^{*}=0$ the algorithm reduces to a forward–backward scheme,
      and the same $O(1/T)$ bound holds.
  \end{itemize}
\end{itemize}

\section*{Discussion}
The theorem shows that the primal–dual hybrid gradient enjoys the optimal
sublinear rate for convex‑concave saddle‑point problems. The proof relies only
on convexity, Lipschitz continuity of $\nabla f$, and the step‑size condition,
hence it applies to a broad class of problems including imaging
inverse‑problems, regularized empirical risk minimization, and constrained
optimization after Lagrangian reformulation.

\newpage
\section*{Preliminaries (Definitions and Lemmas)}
\begin{lemma}[Three‑point identity for the proximal operator]\label{lem:prox}
For any closed convex $h$ and $\tau>0$,
\[
\operatorname{prox}_{\tau h}(u)=\arg\min_{x}\Bigl\{h(x)+\frac{1}{2\tau}\|x-u\|^{2}\Bigr\}
\]
satisfies
\[
\langle x-\operatorname{prox}_{\tau h}(u),\,\operatorname{prox}_{\tau h}(u)-u\rangle
\ge \tau\bigl( h(x)-h(\operatorname{prox}_{\tau h}(u))\bigr),
\qquad\forall x .
\tag{6}
\]
\end{lemma}
\begin{proof}
The optimality condition of the proximal problem yields
$0\in \partial h(\operatorname{prox}_{\tau h}(u)) +\frac{1}{\tau}
(\operatorname{prox}_{\tau h}(u)-u)$. Rearranging gives the inequality \eqref{6}.
\end{proof}

\begin{lemma}[Descent lemma for $L_{f}$‑smooth functions]\label{lem:descent}
If $f$ satisfies (A1), then for any $x,x'$,
\[
f(x')\le f(x)+\langle \nabla f(x),\,x'-x\rangle
      +\frac{L_{f}}{2}\|x'-x\|^{2}.
\tag{7}
\]
\end{lemma}
\begin{proof}
Standard consequence of Lipschitz continuity of the gradient; see e.g.
\cite{Nesterov2004}.
\end{proof}

\section*{Main Proof (Step‑by‑Step Derivation)}
We prove Theorem~\ref{thm:rate}.  
All non‑trivial steps are numbered for easy reference.

\bigskip
\noindent\textbf{Step 1.}  Write the optimality condition for the dual update \eqref{3a}.  
Since $g^{*}$ is convex,
\[
0\in \partial g^{*}(y_{k}) - (Kx_{k}) +\frac{1}{\sigma}(y_{k+1}-y_{k}) .
\tag{8}
\]
Equivalently,
\[
\langle y_{k+1}-y_{k},\,y-y_{k+1}\rangle
\ge \sigma\bigl(g^{*}(y)-g^{*}(y_{k+1}) +\langle Kx_{k},\,y-y_{k+1}\rangle\bigr),
\qquad\forall y .
\tag{9}
\]
(derived from Lemma~\ref{lem:prox} with $h=g^{*}$).

\bigskip
\noindent\textbf{Step 2.}  Write the optimality condition for the primal proximal step \eqref{3b}:
\[
0\in \partial f(x_{k+1}) +\frac{1}{\tau}\bigl(x_{k+1}-(x_{k}-\tau K^{\top}y_{k+1})\bigr).
\tag{10}
\]
Thus,
\[
\langle x_{k+1}-x_{k}+\tau K^{\top}y_{k+1},\,x-x_{k+1}\rangle
\ge \tau\bigl(f(x)-f(x_{k+1})\bigr),
\qquad\forall x .
\tag{11}
\]

\bigskip
\noindent\textbf{Step 3.}  Expand the squared distances that will form our Lyapunov function.
For any $x$,
\[
\begin{aligned}
\|x_{k+1}-x\|^{2}
&=\|x_{k}-x\|^{2}
  -2\langle x_{k+1}-x_{k},\,x_{k+1}-x\rangle
  +\|x_{k+1}-x_{k}\|^{2}.
\end{aligned}
\tag{12}
\]
Similarly for any $y$,
\[
\|y_{k+1}-y\|^{2}
=\|y_{k}-y\|^{2}
 -2\langle y_{k+1}-y_{k},\,y_{k+1}-y\rangle
 +\|y_{k+1}-y_{k}\|^{2}.
\tag{13}
\]

\bigskip
\noindent\textbf{Step 4.}  Multiply \eqref{11} by $1/\tau$ and \eqref{9} by $1/\sigma$, then add the two resulting inequalities.
Using \eqref{12}–\eqref{13} we obtain
\[
\begin{aligned}
&\frac{1}{2\tau}\bigl(\|x_{k}-x\|^{2}-\|x_{k+1}-x\|^{2}\bigr)
 +\frac{1}{2\sigma}\bigl(\|y_{k}-y\|^{2}-\|y_{k+1}-y\|^{2}\bigr) \\
&\qquad
\ge f(x)-f(x_{k+1})
   +g^{*}(y)-g^{*}(y_{k+1})
   +\langle Kx_{k},\,y-y_{k+1}\rangle
   -\langle K^{\top}y_{k+1},\,x-x_{k+1}\rangle \\
&\qquad\qquad
   -\frac{1}{2\tau}\|x_{k+1}-x_{k}\|^{2}
   -\frac{1}{2\sigma}\|y_{k+1}-y_{k}\|^{2}.
\end{aligned}
\tag{14}
\]

\bigskip
\noindent\textbf{Step 5.}  Rearrange the bilinear terms.
Notice that
\[
\langle Kx_{k},\,y-y_{k+1}\rangle
-\langle K^{\top}y_{k+1},\,x-x_{k+1}\rangle
= \langle Kx_{k},y\rangle -\langle Kx_{k},y_{k+1}\rangle
   -\langle Kx_{k+1},y_{k+1}\rangle +\langle Kx_{k+1},y\rangle .
\tag{15}
\]
Insert \eqref{15} into \eqref{14} and collect the saddle‑point terms:
\[
\begin{aligned}
&\frac{1}{2\tau}\bigl(\|x_{k}-x\|^{2}-\|x_{k+1}-x\|^{2}\bigr)
 +\frac{1}{2\sigma}\bigl(\|y_{k}-y\|^{2}-\|y_{k+1}-y\|^{2}\bigr) \\
&\ge
\bigl[f(x_{k+1})+g^{*}(y_{k+1})+\langle Kx_{k+1},y_{k+1}\rangle\big]
-\bigl[f(x)+g^{*}(y)+\langle Kx,y\rangle\bigr] \\
&\quad
+\langle K(x_{k+1}-x_{k}),\,y-y_{k+1}\rangle
-\frac{1}{2\tau}\|x_{k+1}-x_{k}\|^{2}
-\frac{1}{2\sigma}\|y_{k+1}-y_{k}\|^{2}.
\end{aligned}
\tag{16}
\]

\bigskip
\noindent\textbf{Step 6.}  Upper‑bound the cross term.
Using Cauchy–Schwarz and the operator norm of $K$,
\[
\langle K(x_{k+1}-x_{k}),\,y-y_{k+1}\rangle
\le \|K\|\|x_{k+1}-x_{k}\|\|y-y_{k+1}\|
\le \frac{\tau\|K\|^{2}}{2}\|y-y_{k+1}\|^{2}
        +\frac{1}{2\tau}\|x_{k+1}-x_{k}\|^{2}.
\tag{17}
\]
Insert \eqref{17} into \eqref{16}. The $\frac{1}{2\tau}\|x_{k+1}-x_{k}\|^{2}$ terms cancel,
and we obtain
\[
\begin{aligned}
&\frac{1}{2\tau}\bigl(\|x_{k}-x\|^{2}-\|x_{k+1}-x\|^{2}\bigr)
 +\frac{1}{2\sigma}\bigl(\|y_{k}-y\|^{2}-\|y_{k+1}-y\|^{2}\bigr) \\
&\ge
\mathcal L(x_{k+1},y_{k+1})-\mathcal L(x,y)
-\Bigl(\frac{1}{2\sigma}-\frac{\tau\|K\|^{2}}{2}\Bigr)\|y-y_{k+1}\|^{2}
-\frac{1}{2\sigma}\|y_{k+1}-y_{k}\|^{2}.
\end{aligned}
\tag{18}
\]

\bigskip
\noindent\textbf{Step 7.}  Choose $x=x^{\star}$, $y=y^{\star}$ (any saddle point) and drop the last non‑positive term $-\frac{1}{2\sigma}\|y_{k+1}-y_{k}\|^{2}\le0$.
Because of the step‑size condition (A4), we have
\[
\frac{1}{\sigma}-\tau\|K\|^{2}>0 .
\tag{19}
\]
Hence the coefficient in front of $\|y^{\star}-y_{k+1}\|^{2}$ in \eqref{18} is non‑negative and can be omitted, yielding
\[
\frac{1}{2\tau}\bigl(\|x_{k}-x^{\star}\|^{2}-\|x_{k+1}-x^{\star}\|^{2}\bigr)
 +\frac{1}{2\sigma}\bigl(\|y_{k}-y^{\star}\|^{2}-\|y_{k+1}-y^{\star}\|^{2}\bigr)
 \ge \mathcal L(x_{k+1},y_{k+1})-\mathcal L(x^{\star},y^{\star}).
\tag{20}
\]

\bigskip
\noindent\textbf{Step 8.}  Sum \eqref{20} from $k=0$ to $T-1$.
All telescoping terms collapse:
\[
\frac{1}{2\tau}\|x_{0}-x^{\star}\|^{2}
+\frac{1}{2\sigma}\|y_{0}-y^{\star}\|^{2}
\ge \sum_{k=0}^{T-1}\bigl[\mathcal L(x_{k+1},y_{k+1})-\mathcal L(x^{\star},y^{\star})\bigr].
\tag{21}
\]

\bigskip
\noindent\textbf{Step 9.}  By convex–concave structure,
\[
\mathcal L(\bar x_{T},\bar y_{T})-\mathcal L(x^{\star},y^{\star})
\le \frac{1}{T}\sum_{k=0}^{T-1}\bigl[\mathcal L(x_{k+1},y_{k+1})-\mathcal L(x^{\star},y^{\star})\bigr].
\tag{22}
\]
Combine \eqref{21} and \eqref{22} to obtain
\[
\mathcal L(\bar x_{T},\bar y_{T})-\mathcal L(x^{\star},y^{\star})
\le
\frac{1}{2\tau T}\|x_{0}-x^{\star}\|^{2}
+\frac{1}{2\sigma T}\|y_{0}-y^{\star}\|^{2}.
\tag{23}
\]

\bigskip
\noindent\textbf{Step 10.}  The left‑hand side of \eqref{23} is exactly the primal–dual gap
$\mathcal G_{T}$ defined in \eqref{5}. Hence the bound \eqref{5} follows,
establishing the $O(1/T)$ rate. $\square$

\section*{Rate Derivation (With Constants)}
The bound \eqref{5} can be written explicitly as
\[
\mathcal G_{T}
\le
\underbrace{\frac{\|x_{0}-x^{\star}\|^{2}}{2\tau}}_{C_{x}}
\frac{1}{T}
+
\underbrace{\frac{\|y_{0}-y^{\star}\|^{2}}{2\sigma}}_{C_{y}}
\frac{1}{T}.
\]
Thus the constants $C_{x}$ and $C_{y}$ depend linearly on the initial
distances and inversely on the step sizes. Choosing larger $\tau,\sigma$ (subject to \eqref{2}) reduces the constants.

\section*{Special Cases}
\begin{enumerate}
\item \textbf{Pure primal gradient descent.}  
If $g^{*}\equiv0$ (so $y_{k}=0$ for all $k$) and $K=0$, the update reduces to
$x_{k+1}=x_{k}-\tau\nabla f(x_{k})$, and the bound \eqref{5} collapses to the classic
$O(1/T)$ rate for smooth convex minimization.

\item \textbf{Strongly convex $f$.}  
If $f$ is $\mu$‑strongly convex, a similar derivation yields a linear
convergence rate:
\[
\|x_{k}-x^{\star}\|^{2}+\|y_{k}-y^{\star}\|^{2}\le
\bigl(1-\tau\mu\bigr)^{k}
\bigl(\|x_{0}-x^{\star}\|^{2}+\|y_{0}-y^{\star}\|^{2}\bigr).
\]

\item \textbf{Zero primal regularizer ($f\equiv0$).}  
When $f$ vanishes, the primal step becomes $x_{k+1}=x_{k}-\tau K^{\top}y_{k+1}$.
The same analysis holds, and the bound \eqref{5} still guarantees $O(1/T)$
convergence of the dual variable.
\end{enumerate}

\bigskip
\noindent\textbf{References}\\
\begin{thebibliography}{9}
\bibitem{Nesterov2004}
Y. Nesterov, \emph{Introductory Lectures on Convex Optimization}, Kluwer Academic, 2004.
\end{thebibliography}

\end{document}